\section*{Chapitre \ref{ch:b2:retrosynthesis} --- Rétrosynthèse et stratégie de synthèse multiétape}

\begin{solution}{exo:b2:retrosynthesis:1}
\begin{itemize}
  \item 2-phénylpropan-2-ol~: propanone + \ce{PhMgBr}, ou acétophénone + \ce{CH3MgBr}.
  \item Pentan-3-ol~: propanal + \ce{CH3CH2MgBr}.
  \item 1-cyclohexyléthanol~: éthanal + bromure de cyclohexylmagnésium, ou cyclohexanecarbaldéhyde +
        \ce{CH3MgBr}.
  \item 2-méthylbutan-2-ol~: propanone + \ce{CH3CH2MgBr}, ou butanone + \ce{CH3MgBr}.
  \item 1-phényléthanol~: benzaldéhyde + \ce{CH3MgBr}, ou éthanal + \ce{PhMgBr}.
\end{itemize}
\end{solution}

\begin{solution}{exo:b2:retrosynthesis:2}
\omterm{def:b2:retrosynthesis:disconnection}{Synthon} accepteur $\mathrm{Ph\overset{+}{C}H{-}OH}$, d'équivalent le benzaldéhyde~; \omterm{def:b2:retrosynthesis:disconnection}{synthon} donneur \ce{CH3-},
d'équivalent le bromure de méthylmagnésium (ou le méthyllithium). (L'autre coupure, avec \ce{Ph-} comme
donneur, utilise \ce{PhMgBr} et l'éthanal.)
\end{solution}

\begin{solution}{exo:b2:retrosynthesis:3}
$0.90 \times 0.85 \times 0.80 \times 0.95 \times 0.70 = 0.407$~: environ 41~\%.
\end{solution}

\begin{solution}{exo:b2:retrosynthesis:4}
\ce{LiAlH4}, nécessaire pour réduire l'ester en alcool, réduirait aussi la cétone. Protéger la cétone
sous forme d'acétal cyclique (éthane-1,2-diol, catalyseur acide, élimination de l'eau)~; réduire l'ester
par \ce{LiAlH4}~; retirer l'acétal par un acide aqueux dilué~: on obtient la 5-hydroxypentan-2-one.
\end{solution}

\begin{solution}{exo:b2:retrosynthesis:5}
IGF~: butane-1,3-diol $\Rightarrow$ 3-hydroxybutanal (réduction de l'aldéhyde, \ce{NaBH4}). \omterm{def:b2:retrosynthesis:disconnection}{Déconnexion}
à deux groupes de cet aldéhyde $\beta$-hydroxylé au niveau de la liaison $\alpha$--$\beta$~: deux molécules
d'éthanal (\omterm{def:b2:enolates-aldol:aldol}{aldolisation}). Sens direct~: éthanal, base diluée, à froid~; puis \ce{NaBH4}.
\end{solution}

\begin{solution}{exo:b2:retrosynthesis:6}
Le cycle du cyclohexène est coupé au niveau des deux liaisons \ce{C-C} formées dans une \omterm{def:b2:frontier-orbitals:diels-alder}{réaction de
Diels--Alder}, celles qui sont allyliques par rapport à la double liaison du cycle, du côté du
substituant~: buta-1,3-diène et but-3-én-2-one, le groupe acétyle étant porté par le \omterm{def:b2:frontier-orbitals:diels-alder}{diénophile}.
\end{solution}

\begin{solution}{exo:b2:retrosynthesis:7}
En numérotant C1 (\ce{C=O}), C2=C3, puis C4 portant les deux méthyles
(\cref{met:b2:conjugate-additions:robinson-disconnection})~: la coupure aldolique C2--C3 laisse un
aldéhyde en C3 et une méthylcétone~; la coupure de Michael C4--C5 laisse le 2-méthylpropanal (C3 son
carbonyle, C4 son carbone $\alpha$) et la but-3-én-2-one.
\end{solution}

\begin{solution}{exo:b2:retrosynthesis:8}
En une étape~: benzène avec le chlorure de butanoyle et \ce{AlCl3} (\omterm{def:b2:aromatic-substitution:friedel-crafts}{acylation de Friedel--Crafts}, sans
réarrangement, acylation unique). En deux étapes~: benzaldéhyde avec le bromure de propylmagnésium, puis
oxydation de l'alcool secondaire. La voie de Friedel--Crafts est plus courte et ses réactifs moins chers,
mais elle consomme un équivalent entier de chlorure d'aluminium, qui finit dans les effluents aqueux.
\end{solution}

\begin{solution}{exo:b2:retrosynthesis:9}
Le chlorochromate de pyridinium dans le dichlorométhane, l'oxydation de Swern ou le périodinane de
Dess--Martin donnent l'hexanal~; le dichromate ou le permanganate en solution aqueuse acidifiée donnent
l'acide hexanoïque, parce que dans l'eau l'aldéhyde forme son hydrate, qui est oxydé à son tour
(\cref{prop:b2:retrosynthesis:mild-oxidation}).
\end{solution}

\begin{solution}{exo:b2:retrosynthesis:10}
L'amine réagit en premier, non protégée, si l'alcool est protégé~: protéger l'alcool sous forme d'éther
silylé (il résiste à l'acylation)~; acyler l'amine~; retirer le groupe silylé par le fluorure~; acyler
l'alcool. S'il fallait garder l'amine libre pendant que l'alcool réagit, on la protégerait sous forme de
carbamate de \textit{tert}-butoxycarbonyle (retiré par un acide), orthogonal à l'éther silylé (retiré par
le fluorure)~: chacun peut être libéré sans toucher à l'autre.
\end{solution}

\begin{solution}{exo:b2:retrosynthesis:11}
\ce{CH3COCH2CH2COCH3} est un composé 1,4-dicarbonylé~: couper la liaison \ce{C-C} centrale. Le donneur
est l'\omterm{def:b2:enolates-aldol:enolate}{énolate} du 3-oxobutanoate d'éthyle (éthanolate de sodium)~; l'accepteur, un carbone $\alpha$ d'une
cétone, est fourni par la 1-bromopropan-2-one, \ce{BrCH2COCH3} (substitution bimoléculaire du bromure).
Ensuite, l'hydrolyse aqueuse de l'ester et le chauffage éliminent \ce{CO2}~: l'hexane-2,5-dione.
\end{solution}

\begin{solution}{exo:b2:retrosynthesis:12}
Rétrosynthèse~: la méthylcétone flanquée d'un \ce{CH2} est le produit de la synthèse acétylacétique~;
couper la liaison entre C3 et C4~: l'\omterm{def:b2:enolates-aldol:enolate}{énolate} du 3-oxobutanoate d'éthyle et l'halogénure allylique
1-bromo-3-méthylbut-2-ène, \ce{(CH3)2C=CHCH2Br}. Sens direct~: éthanolate de sodium dans l'éthanol, puis
l'halogénure~; puis base aqueuse, acidification et chauffage, qui hydrolysent l'ester et éliminent \ce{CO2}
(\cref{prop:b2:enolates-aldol:decarboxylation})~: la 6-méthylhept-5-én-2-one.
\end{solution}

\begin{solution}{pb:b2:retrosynthesis:1}
\textbf{1.} \ce{HO-C6H4-CH2CH2COCH3}, le \ce{OH} en para de la chaîne~: un phénol et une cétone.
\textbf{2.} L'hydrogénation d'une \ce{C=C}~: \ce{ArCH2CH2CO} $\Rightarrow$ \ce{ArCH=CHCO}, une cétone
$\alpha,\beta$-insaturée.
\textbf{3.} La (\textit{E})-4-(4-hydroxyphényl)but-3-én-2-one, \ce{HOC6H4CH=CHCOCH3}.
\textbf{4.} Un composé carbonylé $\alpha,\beta$-insaturé~: une \omterm{def:b2:enolates-aldol:aldol}{crotonisation}, coupée au niveau de la \ce{C=C}.
\textbf{5.} Le 4-hydroxybenzaldéhyde et la propanone.
\textbf{6.} L'aldéhyde n'a pas d'hydrogène $\alpha$ et ne peut être que l'électrophile~; la propanone, en
excès, est la seule source d'\omterm{def:b2:enolates-aldol:enolate}{énolate} et réagit avec l'aldéhyde, plus réactif, plutôt qu'avec elle-même.
\textbf{7.} (1) 4-hydroxybenzaldéhyde, propanone en excès, hydroxyde de sodium aqueux~; acidification. (2)
\ce{H2} sur palladium sur charbon, à température ambiante, jusqu'à absorption d'un équivalent.
\textbf{8.} Le groupe \ce{CHO} stabilise le phénolate, si bien que ce phénol est au moins aussi acide que le
phénol ($\mathrm pK_a$ 9.99). À $\mathrm{pH} > 13$, $[\ce{ArO-}]/[\ce{ArOH}] > 10^{13-10} = 10^3$~: c'est
le phénolate.
\textbf{9.} Le \ce{O-} repousse la densité électronique à travers le cycle vers le carbone du carbonyle (un
donneur mésomère en para)~: l'aldéhyde est moins électrophile, et la condensation plus lente.
\textbf{10.} Un éther benzylique~: l'hydrogénation sur palladium qui réduit la \ce{C=C} coupe aussi
l'éther benzylique (hydrogénolyse), ce qui libère le phénol dans la même étape.
\textbf{11.} Trois~: benzylation, \omterm{def:b2:enolates-aldol:aldol}{crotonisation}, hydrogénation avec déprotection.
\textbf{12.} Non~: dans des conditions douces (température ambiante, palladium), le cycle \omterm{def:b2:huckel:aromatic}{aromatique} n'est pas réduit.
\textbf{13.} Ne pas protéger~: le phénolate ne fait que ralentir la condensation, ce que compense une
durée plus longue ou un chauffage~; une étape économisée vaut plus que le gain de vitesse.
\textbf{14.} Le 4-iodophénol et la but-3-én-2-one.
\textbf{15.} \ce{HOC6H4I + CH2=CHCOCH3 + Et3N -> HOC6H4CH=CHCOCH3 + Et3NH+ + I-}, catalyseur au palladium.
\textbf{16.} Sur le \ce{CH2} terminal, le carbone $\beta$~: le groupe aryle s'insère à l'extrémité la moins
encombrée, et la $\beta$-élimination d'hydrure donne ensuite l'énone (\textit{E}) conjuguée.
\textbf{17.} \omterm{def:b2:catalytic-cycles:oxidative-addition}{Addition oxydante} de l'iodure d'aryle sur \ce{Pd(0)}, coordination et insertion de l'alcène,
$\beta$-élimination d'hydrure, et régénération de \ce{Pd(0)} par la base, qui capte \ce{HI}
(\cref{ch:b2:catalytic-cycles}).
\textbf{18.} $162.188/(220.005 + 70.091 + 101.193) = 162.188/391.289 \approx 41.4~\%$.
\textbf{19.} \ce{C7H6O2 + C3H6O -> C10H10O2 + H2O}~: $162.188/(122.123 + 58.080) = 162.188/180.203
\approx 90.0~\%$.
\textbf{20.} 100~\%~: chaque atome de \ce{H2} et de l'énone se retrouve dans le produit.
\textbf{21.} $164.204/(122.123 + 58.080 + 2.016) = 164.204/182.219 \approx 90.1~\%$.
\textbf{22.} La \ce{C=C}~: conjuguée et non encombrée, elle se fixe sur le palladium et s'hydrogène bien plus
vite que le \ce{C=O} de la cétone~; s'arrêter après un équivalent de dihydrogène préserve la cétone.
\textbf{23.} Voie aldolique~: aldéhyde, propanone et hydroxyde de sodium bon marché, eau comme sous-produit.
Voie de Heck~: un iodure d'aryle, un catalyseur au palladium et la but-3-én-2-one, très toxique.
\textbf{24.} La voie aldolique~: deux étapes, des réactifs moins chers et moins dangereux, une économie
d'atomes supérieure à 90~\%.
\textbf{25.} Elle le multiplie par 0.90~: $0.782 \times 0.90 \approx 70~\%$.
\textbf{26.} $0.85 \times 0.92 = \boldsymbol{\approx 78~\%}$ au total.
\end{solution}
